% answers_short.tex —— 1999 年·上海卷·理工农医类 · 简短答案速查
% 由 tools/make_short_answers.py 自 answers.tex 的 \gksolhead/\gkanswer 提取，
% 并与 1999/1999_数学真题_上海理_答案要点.md 逐题比对后生成，勿手改（改源 MD 后重跑）。

\begin{gknotice}
\textbf{使用说明}：本节只给最终答案便于快速核对，完整推导见解析版。答案为固定来源 PDF 卷内「答案要点」及解答所载（出版者/版次未核实），非官方标准答案；第 19 题辐角区间按卷内取 $[0,2\pi)$。\textbf{非官方标准答案}。
\end{gknotice}

\gkpart{填空题}
\begin{itemize}[leftmargin=2.6em]
  \item[\textbf{1.}] $2-\sqrt3$。
  \item[\textbf{2.}] $[3,+\infty)$。
  \item[\textbf{3.}] $40$。
  \item[\textbf{4.}] $2\sqrt3$ 或 $\sqrt3$。
  \item[\textbf{5.}] $(2,2)$。
  \item[\textbf{6.}] $\pi$。
  \item[\textbf{7.}] $\left(-\dfrac{5\pi}6,-\dfrac{\pi}3\right)$（源作开区间；函数在闭区间上同样严格递减）。
  \item[\textbf{8.}] $\left\{\dfrac{\sqrt2}{2},\dfrac{3\sqrt2}{2}\right\}$。
  \item[\textbf{9.}] $(\sqrt{10},0)$、$(\sqrt{10},\pi)$。
  \item[\textbf{10.}] $9$。
  \item[\textbf{11.}] $\dfrac29$。
  \item[\textbf{12.}] $\dfrac{\sqrt{11}}{12}$、$\dfrac{\sqrt{14}}{12}$、$\dfrac{\sqrt{11}}6$（三者任一）。
\end{itemize}

\gkpart{选择题}
\begin{center}\begin{tabular}{ccccc}
\toprule
题号 & 13 & 14 & 15 & 16 \\
\midrule
答案 & C & D & B & C \\
\bottomrule
\end{tabular}\end{center}

\gkpart{解答题}
\begin{itemize}[leftmargin=2.6em]
  \item[\textbf{17.}] $0\le a\le1$。
  \item[\textbf{18.}] $\dfrac32\lg2$。
  \item[\textbf{19.}] $\left[0,\dfrac\pi4\right)\cup\left(\dfrac{7\pi}4,2\pi\right)$（卷内辐角区间取 $[0,2\pi)$）。
  \item[\textbf{20.}] （1）$BE\perp PD$；（2）$\arccos\dfrac{\sqrt2}{4}$。
  \item[\textbf{21.}] （1）$\dfrac{2\sqrt6}{3}$ 米；（2）$\dfrac{400\sqrt6}{9}$ 立方米；（3）$\dfrac{\sqrt2}{2}$ 米。
  \item[\textbf{22.}] （1）$\left(\dfrac{a}{\sqrt2},\dfrac{\sqrt2}{a}\right)$（$a>\sqrt2$）；（2）$S(a)$ 的值域为 $(0,\sqrt2)$；（3）$\displaystyle f(a)=\min\{g(a),S(a)\}=\begin{cases}a^2-\dfrac{4}{a^2},&\sqrt2<a\le\sqrt[4]{6},\\[2pt]\sqrt{2\left(1-\dfrac{4}{a^4}\right)},&a>\sqrt[4]{6}.\end{cases}$
\end{itemize}
